\section{The Principles of Well-bred Men}

\textbf{Julius Evola} famously said, in denying he was a fascist, that his principles were those that ``every well-bred man considered healthy, sane, and normal prior to the French revolution.'' It is time, now, to investigate, describe, and defend those principles.

There are two phases in the recovery of Tradition, i.e., the opposite of the regression of the castes:

\begin{enumerate}


\item The first relies on one's own efforts and the development of an inner strength. No longer enthralled by the myths of modernism, man learns to rely on himself. Through self-discipline, he resists the mediocrity of the collective. He integrates his personality into a coherent whole, with all its parts subject to his spirit and his will. No longer dominated by sensuality, greed, willfulness, and attraction to the glamour of the world, he will discover new forces to guide his mind and actions.

\item The first phase is the preparation and the opening to the second. This phase involves an opening to the transcendent and the supernatural.
\end{enumerate}


For our purposes, we will be focusing on the first phase. Evola identified \textbf{Joseph de Maistre} and \textbf{Donoso Cortes} as two of the spirits who exemplified these principles. A contemporary of Evola, \textbf{Charles Maurras} recapitulated that spirit in a natural way, hence we will be using him as a model. This is our promised comparison of those two counter-revolutionaries, in fulfillment of \textbf{Alain de Benoist}'s challenge.


There will be resistance to this, first of all, because the modern mind is totally attached to revolution. Secondly, these two writers are associated with losers, hated losers at that: Evola, through his association with Fascism and Maurras as a collaborator with Vichy France. Yet, those historical accidents do not alter the principles. Evola supported fascism only to the extent it embodied some of those principles. Maurras despised and distrusted the Nazis, but he hated liberalism even more.

Maurras, like Evola, was a classicist, an admirer of the spirit of ancient Greece and Rome. He was also a Medievalist although, again like Evola, he rejected Christianity. Yet there is a fundamental difference between the two; Maurras was a Positivist, totally rejecting the transcendent and the supernatural. Nevertheless, he accepted the reality of the natural law, both for things and for men. Hence, he accepted the moral and social teachings of the Church, particular as exemplified by \textbf{Thomas Aquinas}. He regarded them as rationally demonstrable and in continuity with antique paganism. This actually serves as a benefit, since the principles can be discussed independently of any spiritual or religious system.

Therefore, I shall be translating some of Maurras' works and relating them to those natural law teachings. But first, I need to take a little time off for vacation and translation.

\flrightit{Posted on 2012-05-24 by Cologero}

\begin{center}* * *\end{center}

\begin{footnotesize}
\begin{sffamily}
\texttt{Matt on 2012-05-25 at 15:23 said:}

I've read that Maurras collaborated with a journal that replaced Action Francaise in the final year of his life when he was released from prison into hospital. Are you aware if he did anymore writing at that point (I'm guessing it would be unlikely considering the state he was in)? I know Maurras apparently returned to the Church's metaphysical teachings (and in turn returned to the Church fold) when he was nearing the very end of his life, so the potential development that his writings could have taken from that would be interesting.

\hfill

\texttt{Graham on 2012-05-26 at 12:18 said:}

Regarding the first phase, it's critical to have an idea of what you are to become. An exercise that can help here is to write an obituary for yourself, presuming you became your ideal self before you died.

Without an ideal to guide the individual it's meaningless to think about ``integrating the self'' or ``self-discipline to resist mediocrity''. It is after all our own mediocrity that we must resist, and we can't do so unless we can measure our existing selves against our ideal selves.

\hfill

\end{sffamily}
\end{footnotesize}

